\section{导言：数据不是附属品}

这一讲不是在补一个“数据集名单”，而是在回答一个更根本的问题：\textbf{语言模型到底靠什么长出来。} 前几讲已经讲完了架构、优化器、tokenizer、并行和 scaling laws，默认前提都是“数据已给定”。这一讲把前提翻过来，讨论数据从哪里来、怎么过滤、怎么混合、怎么做成训练集，以及哪些法律和合规边界不能绕过。

\begin{importantbox}{这一讲的核心判断}
\begin{enumerate}
    \item 数据不是训练流水线里的一个附件，而是决定模型上限的主要变量。
    \item 数据工作是长期、并行、可拆分的工程，不是一次性的抓取任务。
    \item 好模型通常不是因为“多看了互联网”，而是因为“更会挑数据”。
    \item 数据处理不仅是技术问题，也是版权、隐私、平台规则与供应链问题。
\end{enumerate}
\end{importantbox}

讲者对”data is the most important thing”这个判断非常强。理由并不抽象：架构往往由少数人定义，数据工作却可以拆成大规模并行流水线，持续调优、持续清洗、持续重配比。模型越大，数据质量、覆盖面和过滤策略就越直接地决定算力是否用在刀刃上。

\begin{importantbox}{Data is the New Moat}
在大模型竞争中，架构设计趋同（几乎所有主流模型都是 decoder-only Transformer），训练框架也日趋标准化。真正构成差异化壁垒的，是\textbf{数据}：你能拿到什么数据、怎么清洗、怎么配比、怎么持续迭代。这就是为什么 Llama 3 的技术报告详细描述了架构和训练流程，却对数据只给出模糊描述——数据配方本身就是最核心的竞争力。对创业公司而言，如果没有独特的数据获取和处理能力，仅靠公开权重或公开代码很难建立持久优势。
\end{importantbox}

\subsection{为什么公司对数据这么保密}

公开模型论文通常会详细讲结构、训练流程和损失函数，但对数据往往只给粗粒度描述。原因有两个：
\begin{itemize}
    \item \textbf{竞争原因}：数据配方本身就是产品竞争力。
    \item \textbf{法律风险}：版权、隐私和条款违约都可能引发争议。
\end{itemize}

这也解释了为什么”公开权重”并不等于”公开数据”。前者解决模型可用性，后者涉及数据供应链和合规问题。

\begin{table}[H]
\centering
\begin{tabular}{p{0.2\textwidth}p{0.35\textwidth}p{0.35\textwidth}}
\toprule
\textbf{公开程度} & \textbf{典型做法} & \textbf{背后的考量} \\
\midrule
公开权重 & 发布 checkpoint，任何人可下载推理 & 推广生态、吸引开发者 \\
公开架构 & 论文详细描述层数、注意力头、激活函数 & 学术声誉、可复现性 \\
数据常隐藏 & 只给粗粒度描述（如”网页+书籍”） & 竞争壁垒 + 法律风险规避 \\
\bottomrule
\end{tabular}
\caption{公开论文里最透明和最不透明的，往往不是同一件事}
\end{table}

\subsection{数据也是长尾问题}

\begin{knowledgebox}{为什么数据是 long-tail problem}
你可以雇很多人、开很多爬虫、写很多过滤规则，但世界上真正稀有、边缘、昂贵、合规复杂的数据，总是以长尾方式分布。架构是少数设计点，数据却是无数细粒度选择的叠加。
\end{knowledgebox}

当你要做一个既懂代码、又懂数学、又会对话、还要能长上下文和安全拒答的模型时，没有哪一种单独数据源可以包打天下。数据质量与覆盖面之间存在天然张力：高质量来源（如 Wikipedia、教科书）通常覆盖面窄；广覆盖的来源（如随机网页）噪声极大。成熟的数据工程，本质上就是在这条光谱上持续做折中。

